\documentclass[numbers=enddot]{scrartcl}

\usepackage{amsmath}
\usepackage{amssymb}
\usepackage{amsthm}
\usepackage[dvipsnames]{xcolor}
\usepackage{stmaryrd}
\usepackage{tikz-cd}

\input{lambdamacs}
% \newcommand{\?}{{.\ }}
%\theoremheaderfont{\scshape}
%\theorembodyfont{\normalfont}
% \theoremseparator{.\ \ }
\newtheorem{thm}{Theorem}
\newtheorem{prop}[thm]{Proposition}
\newtheorem{cor}[thm]{Corollary}
\newtheorem{lemma}[thm]{Lemma}
\theoremstyle{definition}
\newtheorem{rmk}[thm]{Remark}
\newtheorem{dfn}[thm]{Definition}
\newtheorem{notn}[thm]{Notation}
\newtheorem{expl}[thm]{Example}
\newtheorem{asmn}[thm]{Assumption}
%\theoremstyle{nonumberplain}
%\theoremsymbol{\Box}


% \usepackage{MnSymbol}

\usepackage{bussproofs}
\EnableBpAbbreviations

\title{Isomorphisms between algebraic datatypes}
\author{Andrew Polonsky, Ben Lenox}

\newcommand{\cat}{\mathbb{C}}
\newcommand{\op}{{op}}
\newcommand{\ints}{\mathbb{Z}}
\newcommand{\mcS}{\mathcal{S}}
\newcommand{\semof}[1]{{\left\llbracket{#1}\right\rrbracket}}
\newcommand{\cc}[1]{{[#1]_\circlearrowleft}}
\newcommand{\downg}[1]{{#1}{\downarrow_G}}
\newcommand{\Xcyc}{X_{\circlearrowright}}
\newcommand{\Xacyc}{X_{\not\circlearrowright}}
\newcommand{\downe}[1]{{#1}{\downarrow_E}}
\newcommand{\LRA}{\Longrightarrow}
% \newcommand{\mcS}{\mathpzc{S}}

\newcommand{\nullX}{\mathpzc{N}}
\newcommand{\gtE}{\geq_E}
\newcommand{\Xmon}[1]{X^{\scriptscriptstyle{\ge}}_{#1}}
\newcommand{\cov}{\sqsupseteq}
\newcommand{\LA}{\Leftarrow}

 \newcommand{\mzero}{\ul{\mathbf{0}}}
 \newcommand{\mone}{\ul{\mathbf{1}}}
 \newcommand{\const}[1]{\ul{\mathbf{#1}}}
\begin{document}

\tableofcontents

\section{Definitions and Notations}

Fix a finite set of variables $X$.  %We write $X = \setof{x_1,\dots,x_n}$.

For each $x \in X$, let there be given a polynomial $p_x \in \nat[X]$.

We consider the system $E$ of equations $\setof{x = p_x \mid x \in X}$.
To stress the dependence of $p_x$ on the variables in $X$, we may write
$p_x$ as $p_x(X)$, $p_x(x_1,\dots,x_n)$, or $p_x(\bar x)$, with the
understanding that $X = \setof{x_1,\dots,x_n}$.

We consider polynomials up to semiring axioms.
Any such polynomial is a sum of monomials over variables from $X$ with non-zero coefficients.
For our purposes however, it is more convenient to consider each monomial as a pure
 products of variables.

 That is, we write each polynomial $p$ as $p = \sum m_i$, where each monomial $m_i$
 has coefficient 1, but may occur more than once in this sum.
 Similarly, each monomial $m$ is represented as $m = \prod x_i$
 where multiple $x_i$s may denote the same element of $X$.
 % A constant polynomial $n$ is represented as $\sum_{i = 1}^n m$, where
 % $m = \prod_{i \in \emptyset} x_i$.

 \begin{notn}\hfill
   \begin{itemize}
     \item $\mone = \prod_{\emptyset}$ is the monomial representing 1.
     \item $\mzero = \sum_\emptyset$ is the polynomial representing 0.
     \item $\const{n} = \sum_{i = 1}^n \mone$ is the polynomial representing $n > 0$.
   \end{itemize}
 \end{notn}

\section{Setup}

\subsection{Equality modulo $E$}
\hired{Need to reframe this discussion.
Removing acyclic stuff is completely sound for the semiring or rig category setting,
doesn't change the ri(n)g/groupoid.
But removing the nullable stuff does change it.
Need to motivate it, and have two points to make:
The primary focus is ADT isomorphisms, for which it IS sound,
and second is that ``removing the nullables" is an algebraically
universal/structural operation, and it does reveal something about the free rig as well,
with the connection being made explicit by a quotient map.
Conjecture: on the high elements, this map admits a section,
an ``inclusion" of $\nat[X_{\geq 1}]/E_{\geq 1} \to \nat[X]/E$.
}
The principal question we will investigate in this work is under which conditions
the entailment $E \models p = q$ holds in the theory of semirings, where
$p$ and $q$ are arbitrary polynomials over $X$.

\begin{notn} \hfill
  \begin{itemize}

  \item $p=q$ denotes equality of polynomials $p$ and $q$ as elements of
  $\nat[X]$.\\
  That is, $p = q$ is equality up to rig axioms.

  \item $p=_E q$ is the assertion that for any rig
  $\mcS$ that models $E$, $\semof{p}^\mcS = \semof{q}^\mcS$.\\
  That is, $p =_E q$ is entailment $E \models_{Rig} p=q$.
\end{itemize}
\end{notn}

Elementary universal algebra yields the standard characterization:
$p =_E q$ iff $\semof{s} = \semof{t}$ in the free semiring over $X$,
quotiented by the congruence generated by $E$.  Thus, the question reduces to the
study of equality in the semiring $\nat[X]/E$. Given the special format of
equations in $E$, we will later introduce a rewriting system precisely for this purpose.

\subsection{High elements of the rig $\nat[X]/E$}
Our next goal is characterize the high elements of $\nat[X]$: polynomials
$p$ such that for all $q$, $p =_E q + r$ for some $r$.

\hired{Discuss Fiore--Leinster conditions: $x \ge x^2$, $x \ge 1$, $2x \ge x$.\\
In fact, it suffices to prove $x \ge x^2$, $x \ge 2$.}

Generalizing these conditions to the multivariate case
requires a more refined analysis of $E$-equality, for which we introduce a rewrite system
that we show to be confluent.
% Roughly, these will be polynomials containing a monomial that can generate all other variables.
% Moreover, each variable $x$ in this monomial should satisfy $x \gtE 2$ and $x \gtE x^2$.

\subsection{Unfolding reductions}
\begin{dfn} \hfill
  \begin{enumerate}
\item
  The \emph{unfolding reduction} is the rewrite rule
  \[ x \rrule p_x \]
\item  This notion is extended to arbitrary polynomials in an obvious way:
  $q \lra q'$ iff $q'$ results from $q$ by replacing a single occurrence of
  a variable $x$ with $p_x$.
  \item We write $q \LRA q'$ if $q'$ results from $q$ by replacing zero or
   more variable occurrences with the corresponding polynomials ---
   for example, replacing some occurrences of $x$ with $p_x$, some occurrences of $y$ with $p_y$, etc. --- in one step.
   \hired{Formally, the rules generating $\LRA$ are as follows:
   $x \LRA x$, $x \LRA p_x$, and congruence under $\Pi$ and $\Sigma$.}
\item
  We write $q \thra q'$ if $q'$ results from $q$ after zero or more unfolding
  reductions starting from $q$. That is, $\thra$ is a reflexive--transitive closure of $\lra$.
  \item $\downe{p} = \setof{q \mid p \thra q}$
\end{enumerate}
\end{dfn}

\begin{rmk}
The relations $\lra, \LRA, \thra$ above are defined
    for \emph{polynomials}; since these are represented as sums of monomials,
    there is implicit conversion of the source and target of any unfolding reduction
    to such representations.

    In particular, we do \emph{not} have $x(y + z) \lra p_x(y + z)$.

    In contrast, we \emph{do} have $x(y + z) \LRA p_x(y + z)$ and $x(y + z) \thra p_x(y + z)$.
\end{rmk}

Our next goal is to prove a confluence theorem for unfolding reductions.
This follows by a straightforward adaptation of the usual Martin-L\"of argument.

\begin{prop} \hfill
    \label{p:lra}
  \begin{enumerate}
    \item For $R \in \setof{\lra,\LRA,\thra}$, $pRq$ implies $\downe{q} \subseteq \downe{p}$.
    \item $\lra \subseteq \LRA \subseteq \thra$. \label{p:lra-LRA}
    % \item $\thra$ is transitive, hence
    \item $\LRA^* \subseteq \thra$, where $\LRA^*$ is the reflexive--transitive closure of $\LRA$.
    \label{p:lra-thra}
    \item $\LRA$ is a congruence with respect to addition and multiplication, in particular
    \item \label{p:lra-cong} $p \LRA p'$ implies $p q \LRA p' q$.
  \end{enumerate}
\end{prop}
\begin{proof}
Routine.
\end{proof}

\begin{lemma}(DP and PML for unfolding reductions)
    \label{l:PML}
  \begin{enumerate}
    \item The parallel unfolding reduction satisfies the strong diamond property:
    \[\begin{tikzcd}
      {s} & {u} \\
      {t} & {v}
      \arrow[Rightarrow, from=1-1, to=1-2]
      \arrow[Rightarrow, from=1-1, to=2-1]
      \arrow[Rightarrow, dashed, from=1-2, to=2-2]
      \arrow[Rightarrow, dashed, from=2-1, to=2-2]
    \end{tikzcd}
    \]
    \item If $s \LRA t$ and $s \thra u$ then there exists a $v$ such that $t \thra v$ and $u \LRA v$.
    \[\begin{tikzcd}
      {s} & {u} \\
      {t} & {v}
      \arrow[two heads, from=1-1, to=1-2]
      \arrow[Rightarrow, from=1-1, to=2-1]
      \arrow[Rightarrow, dashed, from=1-2, to=2-2]
      \arrow[two heads, dashed, from=2-1, to=2-2]
    \end{tikzcd}
    \]
    \item If $s \lra t$ and $s \thra u$ then there exists a $v$ such that $t \thra v$ and $u \thra v$.
  \[\begin{tikzcd}
    {s} & {u} \\
    {t} & {v}
    \arrow[two heads, from=1-1, to=1-2]
    \arrow[from=1-1, to=2-1]
    \arrow[two heads, dashed, from=1-2, to=2-2]
    \arrow[two heads, dashed, from=2-1, to=2-2]
  \end{tikzcd}
  \]
\end{enumerate}
\end{lemma}

\begin{proof}
  \begin{enumerate}
    \item
    Recall that a general polynomial $s$ can be written as $s = \sum s_i$,
    with each $s_i$ a product of variables (any monomial may appear more than once in this sum).

    Since parallel reduction is a congruence with respect to addition, it suffices
    to verify this statement when $s$ is such a monomial: $s = \prod_{i \in I} x_i$, with each $x_i \in X$.

    \hired{Not true, we really need to define $\LRA$ as the congruence wrt $\Pi$ and $\Sigma$.}

    Let $T$ be the set of $i \in I$ for which $x_i$ is unfolded in $s \LRA t$.

    Let $U$ be the set of $i \in I$ for which $x_i$ is unfolded in $s \LRA u$.

    Then $t = \prod_{i \in T} p_{x_i} \times \prod_{i \in I \setminus T} x_i$.

    Also $u = \prod_{i \in U} p_{x_i} \times \prod_{i \in I \setminus U} x_i$.

    Put $v = \prod_{i \in T \cup U} p_{x_i} \times \prod_{i \in I \setminus (T \cup U)} x_i$.

    Then $t \LRA v$ by applying Proposition \ref{p:lra}.\ref{p:lra-cong} with $p = \prod_{i \in U \setminus T} x_i$.

    Similarly, we get $u \LRA v$.

    \item By tiling the diagram below using 1, and applying Proposition \ref{p:lra}.\ref{p:lra-thra}
    to the bottom edge of the diagram.
  \[ \begin{tikzcd}
  	{s=s_0} & {s_1} & {s_2} & {\cdots} & {s_n=u} \\
  	{t=t_0} & {t_1} & {t_2} & {\cdots} & {t_n=v}
  	\arrow[rightarrow, from=1-1, to=1-2]
  	\arrow[Rightarrow, from=1-1, to=2-1]
  	\arrow[rightarrow, from=1-2, to=1-3]
  	\arrow[Rightarrow, dashed, from=1-2, to=2-2]
  	\arrow[rightarrow, from=1-3, to=1-4]
  	\arrow[Rightarrow, dashed, from=1-3, to=2-3]
  	\arrow[equals, dotted, from=1-4, to=1-5]
  	\arrow[Rightarrow, dashed, from=1-4, to=2-4]
  	\arrow[Rightarrow, dashed, from=1-5, to=2-5]
  	\arrow[Rightarrow, dashed, from=2-1, to=2-2]
  	\arrow[Rightarrow, dashed, from=2-2, to=2-3]
  	\arrow[Rightarrow, dashed, from=2-3, to=2-4]
  	\arrow[equals, dotted, from=2-4, to=2-5]
  \end{tikzcd} \]
  \item By the previous point and by parts \ref{p:lra-LRA} and \ref{p:lra-thra} of Proposition \ref{p:lra}.
  \qedhere
    \end{enumerate}
\end{proof}

\begin{prop}
  \label{p:CR}
The unfolding reduction is confluent modulo semiring axioms:
\[ p =_E q \iff \downe{p} \cap \downe{q} \neq \emptyset \]
\end{prop}

\begin{proof}
The right-to-left implication is trivial.

To prove the left-to-right implication, we proceed by induction on $p=_E q$.
Since $=_E$ is the least congruence containing $\lra$ and semiring equality,
any proof that $p=_E q$ is a finite sequence of $\lra$ reduction or expansion steps, interleaved with
applications of semiring axioms.  We proceed by induction on the length of this sequence.
\begin{description}
\item[Base case.] If $p = q$ by the semiring axioms, then $\downe{p} = \downe{q}$, and both contain $p$.
\item[Induction 1.] If $p \lra p' =_E q$, then by induction hypothesis,
there exists $r \in \downe{p'} \cap \downe{q}$.  Since $\downe{p'} \subseteq \downe{p}$,
also $r \in \downe{p} \cap \downe{q}$.
\item[Induction 2.] If $p \longleftarrow p' =_E q$ then we proceed as follows.

By IH, let $r \in \downe{p'} \cap \downe{q}$, so that $p' \thra r$ and $q \thra r$.

By Lemma \ref{l:PML}, there exists a $u$ such that
$p \thra u$ and $r \thra u$.

So $q \thra r \thra u$ and $u \in \downe{p} \cap \downe{q}$. \qedhere
\end{description}
\end{proof}

\begin{prop} Let $m_i$ be monomials and $q$ be a polynomial.
  \label{p:sum-red}
  \begin{enumerate}
    \item
    $\sum m_i \LRA q \iff q = \sum q_i \text{ and }\ \forall i. m_i \LRA q_i$
    \item
    $\sum m_i \thra q \iff q = \sum q_i \text{ and }\ \forall i. m_i \thra q_i$
  \end{enumerate}
\end{prop}
\begin{proof}
  \begin{enumerate}
    \item
    This is just the definition of $\LRA$ for polynomials.
    \item
    By induction on the length of the reduction $\Sigma m_i \thra q$. \qedhere
  \end{enumerate}
\end{proof}


\section{The covering preorder}

\begin{dfn} Let $p, q \in \nat[X]$.
\begin{enumerate}
  \item We write $p \gtE q$ if there exists $r \in \nat[X]$
  such that $p =_E q + r$.
  \item We write $p \cov q$ if there exists a $r\in \nat[X]$ such that
  $p \thra q + r$.
  % \item For $p \in \nat[X]$, let $X_{\ge p} = \setof{x \in X \mid x \gtE p}$.
  % \item Let $X_\square = \setof {x \mid x \ge x^2}$.
\end{enumerate}
\end{dfn}

Note that $\gtE$ and $\cov$ are reflexive, transitive, and are congruences
with respect to addition and multiplication.

Trivially, $p \cov q$ implies $p \gtE q$.
Our next goal is to isolate conditions under which
the converse holds as well, so that $p \cov q \iff p \gtE q$.

First we record an auxiliary lemma about the covering relation.

\begin{prop}
  \label{p:cov2} Let $m$ be a monomial.
  \[\prod x_i \cov m \iff m = \prod m_i \text{ and }\ \forall i. x_i \cov m_i\]
\end{prop}

\begin{proof}
The converse is trivial because $\cov$ is a congruence.

The forward direction is proved by induction on $\prod x_i \thra m + q$.

The base case is only possible when $q = 0$ and $m = \prod x_i$. Then $m_i = x_i$ and $x_i \cov m_i$ by reflexivity.

We now consider the inductive step. Suppose $\prod x_i \lra p \thra m + q$.

Then $p = p_{x_j} \times \prod_{i \neq j} x_i$ for some $j$.

To ease the notation, let us write $m_{\neq j} = \prod_{i \neq j} x_i$.

By Proposition \ref{p:sum-red}, there exists a monomial $m'$ in $p$,
so that $p = m' + p'$, and $m' \thra m + q'$.

Applying the induction hypothesis to $m' \thra m + q'$ we find that,
\begin{equation}
   \label{e:yj}
   m = \prod_{y_i \in m'} m_i, \qquad y_i \cov m_i
\end{equation}
As any monomial of $p$, $m'$ has the form $m'' \cdot m_{\neq j}$
for some monomial $m''$ of $p_{x_j}$ --- and so $x_j \cov m''$.

With $m = m'' \times m_{\neq j}$, \eqref{e:yj} can be further refined into
\begin{align*}
   % \label{e:zj1}
   m &= \left(\prod_{i \neq j} m_i\right)
 \cdot \left(\prod_{z \in m''} m_z\right)  \\
 % \label{e:zj2}
   &y_i \cov m_i \qquad (i \neq j) \\
   &m'' \cov m_j := \prod_{z \in m''} m_z
\end{align*}

So we have $y_i \cov m_i$ for all $y_i \in m_{\neq j}$, and
we also have $x_j \cov m'' \cov m_j$ and hence $x_j \cov m_j$.

We conclude that $x_i \cov m_i$ for all $i$, where $m = \prod m_i$. \qedhere
\end{proof}

\subsection{The recurrent elements}

\begin{prop}
  % \begin{enumerate}
    % \item A polynomial $p = \sum m_i$ is recurrent iff each $m_i$ is.
  % \end{enumerate}
 % A monomial $m = \prod x_i$ is recurrent iff $x_i \in \Xsc$ for each $i$.
 Let $m = \prod_{i \in I} x_i$ and suppose each $p_{x_i} \cov x_i$ for all $i$.
 \begin{enumerate}
   \item If $m \LRA q$, then $q \cov m$.
   \item If $m \thra q$, then $q \cov m$.
 \end{enumerate}
\end{prop}

\begin{proof} \hfill
  \begin{enumerate}
  \item Suppose $m \LRA q$.

  Then $q = \big(\prod_{i \in U} p_{x_i} \big) \cdot \big(\prod_{i \in I \setminus U} x_i\big)$ for some $U \sse I$.

  By hypothesis, $p_{x_i} \thra x_i + r_i$ for all $i$.

  Now by Proposition \ref{p:lra}.\ref{p:lra-cong},
  \begin{align*}
    q &\thra \prod_{i \in U} (x_i + r_i) \cdot \prod_{i \in I \setminus U} x_i \\
    &= \left(\prod_{i \in U} x_i  \cdot \prod_{i \in I \setminus U} x_i\right) + r' \\
    &= m + r'
  \end{align*}

  \item
  Suppose $m \thra q$. We show that $q \cov m$ by induction on the length of this reduction.

\begin{description}
  \item[Base case.]
  If $q = m$, then trivially $q \cov m$.
  \item[Induction.]
  Suppose  $m \thra q' \lra q$.  By induction hypothesis, $q' \cov m$.

  So $q' \thra m + r$ for some $r$.

  By the second part of Lemma \ref{l:PML}, there exists $t$ such that
   $m + r \LRA t$ and $q \thra t$.

   Then $t = t_1 + t_2$, where $m \LRA t_1$ and $r \LRA t_2$.

   By the part 1, $t_1 \thra m + r'$.

   Now $q \thra t_1 + t_2 \thra m + r' + t_2$ by congruence.  Indeed, $q \cov m$. \qedhere
\end{description}
\end{enumerate}
\end{proof}

\begin{dfn}
  $p \in \nat[X]$ is \emph{recurrent} iff
  \[ p \thra q \implies q \cov p \]
\end{dfn}

\newcommand{\Xsc}{\Xcyc^+}
Let $\Xsc = \setof{x \in X \mid x \text{ is recurrent}}$.

\begin{cor} \label{p:sc=cov} \hfill
  \begin{enumerate}
    \item
    $x \in \Xsc \iff p_x \cov x$
    \item
    $m = \prod {x_i}$ is recurrent whenever each $x_i$ is.
    \item
    If $x \in \Xsc$ for each variable $x$ occuring in $p$, then $p$ is recurrent.
  \end{enumerate}
\end{cor}

\begin{proof} \hfill
  \begin{enumerate}
    \item $x \in \Xsc$ implies $p_x \cov x$ because $x \lra p_x$.

  The converse holds by the above proposition with $m = x$.

  \item By part 1 and the above proposition.

  \item Write $p = \sum m_i$. If $p \thra q$, then $q = \sum q_i$ with
  $m_i \thra q_i$ for each $i$.

  By part 2, $m_i$ is recurrent, so $q_i \cov m_i$, and
  $q_i \thra m_i + r_i$.

  Now $q \thra \sum(m_i + r_i) = (\sum m_i) + (\sum r_i) = p + \sum r_i$, and $q \cov p$.\qedhere
\end{enumerate}
\end{proof}

\begin{prop}
  \label{p:gtE-cov}
  If $p \gtE q$ and $q$ is recurrent, then $p \cov q$.
\end{prop}
\begin{proof}
  Suppose $p =_E q + r$.

  By Proposition \ref{p:CR}, let $p'$ be such that
  $p \thra p'$ and $q + r \thra p'$.

  By Proposition \ref{p:sum-red}, $p' = q' + r'$ where
  $q \thra q'$ and $r \thra r'$.

  Since $q$ is recurrent, $q' \cov q$, whence
  $q' \thra q + r''$.

  Thus $p \thra p' = q' + r' \thra q + r'' + r'$ and $p \cov q$.
\end{proof}

\subsection{Inductive characterization of variables covering a monomial}
\begin{dfn}
  For any monomial $m$, we inductively define a set $\Xmon{m} \sse X$
  that characterizes those variables $x$ that can generate the monomial $m$ by unfolding reductions.
  % of ``variables that can generate an $m$ by unfolding reductions",
  % which we will later show to be equal to $\setof{x \mid x \cov m}$.

  \[
  \AXC{$\phantom{\Xmon{m_j}}$}
  \RL{Var}
  \UIC{$x \in \Xmon{x}$}
  \DP
  \qquad \qquad
  \AXC{$p_x = p' + \Pi y_j$}
  \AXC{$\forall j.\ y_j \in \Xmon{m_j}$}
  \RL{Mon}
  \BIC{$x \in \Xmon{\Pi m_j}$}
  \DP
  \]
\end{dfn}

\begin{prop}
  \label{p:cov1}
  $x \cov m \iff x \in \Xmon{m}$
\end{prop}

\begin{proof}
   The reverse direction is trivial.

  The forward direction of is proved by induction
  on the length of the reduction to $m + q$.

  In the base case, an empty reduction from $x$ to $m + q$
  yields $q = 0$ and $m = x$, whence $x \in \Xmon{m}$.

  So it remains to consider the inductive step.

    Suppose $x \lra p_x \thra m + q$.

    By Proposition \ref{p:sum-red}, $p_x = p' + m'$, where $m'$ is a monomial,
    $m' \thra m + r$, $p' \thra q'$, and $q = q' + r$.

    Write $m' = \prod y_i$.  % \hired{Then $m' \cov m$ by a shorter reduction than $p_x \cov m$.}
    By Proposition \ref{p:cov2}, $m = \prod m_i$, with each $y_i \cov m_i$.

    These reductions are shorter because they do not include the original unfolding
    $x \lra p_x$; by induction hypothesis, we obtain $y_i \in \Xmon{m_i}$.

    Since $m = \prod m_i$, applying the Mon rule yields $x \in \Xmon{m}$.
\end{proof}

\begin{cor} \label{c:recXmon}
  If $m$ is recurrent, then $x \gtE m \iff x \in \Xmon{m}$.
\end{cor}


\subsection{The compactible elements}

The inductive definition of variables covering a monomial enables us to
effectively decide when a recurrent variable satisfies $x \gtE x^2$.

\begin{prop} \label{p:compact-n}
If $x \gtE x^2$, then $x \gtE x^n$ for any $n$.
\end{prop}
\begin{proof}
    Induction on $n$.
\end{proof}


\begin{prop}%(An effective algorithm for deciding $x \ge x^2$ for recurrent $x$)
  Assume $x$ is recurrent.
The following yields an effective way to decide whether $x \ge x^2$.

By Corollary \ref{p:sc=cov}.2, if $x$ is recurrent, then so is $x^2$.

By the previous corollary, $x \ge x^2$ is equivalent to $x \in \Xmon{x^2}$.

It therefore suffices to decide the latter, and this can be achieved as follows.

\begin{enumerate}
  \item Enumerate all elements of $\Xmon{1}$.
  \begin{verbatim}
    Initialize S = {}.
    For each variable y,
      if py contains a monomial each of whose variables are in S,
        add y to S.
    Repeat while new elements are being added to S.
  \end{verbatim}
    \item Enumerate all elements of $\Xmon{x}$.
  This can be obtained by the following loop.
  \begin{verbatim}
    Initialize S = {x}.
    For each variable y,
      if py contains a monomial m = v * m', m' subset X1, v in S,
        add y to S.
    Repeat while new elements are being added to S.
  \end{verbatim}

  \item  Enumerate all elements of $\Xmon{x \cdot x}$.

  This can be obtained by the following loop.
  \begin{verbatim}
    Initialize S = {}.
    For each variable y,
      if py contains a monomial m = v1 * v2 * m', m' subset X1, v1, v2, in Xx,
                  or a monomial m = v * m', m' subset X1, v in S,
        add y to S.
    Repeat while new elements are being added to S.
  \end{verbatim}
\end{enumerate}
\end{prop}

\section{Characterizing the high elements}


\subsection{The nullable elements}

A key obstruction to the analysis of the rig $\nat[X]/E$ are variables that can never generate a constant.
When the system $E$ is considered as a datatype definition, these variables correspond to empty types.
\begin{expl}
Consider the following system
\begin{align*}
  x &= 1 + y \\
  y &= y + xy
\end{align*}
In this example, $y \not \gtE 1$. As a result, $x \notin \Xsc$.  This causes Proposition
\ref{p:gtE-cov} to fail when the polynomial $q$ contains $x$. Indeed,
\[ x^2 =_E x(1 + y) =_E x + xy =_E 1 + y + xy =_E 1 + y =_E x \]
and thus $x \gtE x^2$.
But $x \not \cov x^2$, for any non-empty reduction from $x$ ends in a term containing $y$, and this $y$ can never be erased by unfolding reductions.

As a datatype definition, this system results in $x$ being the unit type and $y$ being the empty type.
\end{expl}

Our first step is therefore to characterize such elements.

\begin{dfn}
  $x \in X$ is \emph{nullable} if there does \emph{not} exist $q$ such that
  $x =_E q + 1$.

  That is, $x$ is nullable if $x \not \gtE 1$.

  We denote by $\nullX$ the set of all nullable elements of $X$.
\end{dfn}

\begin{prop}
  $x \gtE \const 1 \implies x \cov \const 1$.
\end{prop}
\begin{proof}
  Since $\const 1$ is a normal form for unfolding reductions, it is vacuously recurrent.
  Thus Proposition \ref{p:gtE-cov} applies.
\end{proof}

% \begin{prop}
%   \begin{enumerate}
%     \item
%     If $p + \mone \thra q$, then $q = q' + \mone$.
%     \item
%     If $x =_E p + 1$, then there is $q$ such that $q + 1 \in \downe{x}$.
%   \end{enumerate}
% \end{prop}
% \begin{proof}
%   \begin{enumerate}
%     \item Since unfolding reductions can only be applied to monomials that contain at least one variable.
%     \item By combining part 1 with Proposition \ref{p:CR}. \qedhere
%   \end{enumerate}
% \end{proof}

\begin{prop} \label{p:xcovn}
  For each $x \in \Xsc$, $x \cov \const 1$ implies $x \cov \const{n}$ for each $n$.
\end{prop}
\begin{proof} By induction on $n$, it suffices to show that $x \cov 1 + x$.

  From $x \cov \const{1}$, we have $x \thra 1 + q$ for some $q$.
 
  From $x \in \Xsc$, we have $1 + q \thra x + r$ for some $r$.

  Since $\const{1}$ is a normal form, $r = 1 + r'$ and $q \thra x + r'$.

  So $q \cov x$ and $x \cov \const{1} + q \cov \const{1} + x$
\end{proof}

% Note that the (non-)nullable elements are invariant under $E$-equality, hence, also under unfolding reductions.

% \begin{rmk}
When the system $E$ is used to define a set of mutually recursive datatypes,
all of the nullable variables denote empty types.  That is, the initial algebra
for the functor $F_E(X_1,\dots,X_n) = (p_1(\bar X), \dots,p_n(\bar X)) :
Set^n \to Set^n$ consists of the types $(A_1,\dots,A_n)$, and
$A_i$ is empty iff $x_i \in \nullX$.

The remaining datatypes $A_j$ for $x_j \notin \nullX$ can be described
as fixed points of a functor that does not contain nullable variables at all.

(Namely, the system $E'$ obtained from $E$ by removing the nullable variables.)
% \end{rmk}

In light of the above, in the study of isomorphisms between ADTs,
it suffices to restrict our attention to systems containing only non-nullable
variables. We thus admit the following assumption for the remainder of the paper.

\begin{asmn}
  \label{asmn:null}
   \fbox{$\nullX = \emptyset$}
 \end{asmn}

\hired{Note that this does change the ring,
so it's not "innocuous" as removing acyclic variables.}

\begin{lemma}\label{l:sse-mon-cov} Suppose Assumption \ref{asmn:null} is valid.
    Let $m = \Pi_{i \in I}v_i$.
    Let $J \sse I$.  Then $m \cov \Pi_{j \in J} v_j$.
\end{lemma}
\begin{proof}
    Let $m = \Pi_{i \in I}v_i$ and $J \sse I$.  By Assumption \ref{asmn:null}, for each $i \in I \setminus J$, $v_i \gtE 1$.  Because $\cov$ preserves multiplication, 
    \[ m = \prod_{i \in I}v_i \cov (\prod_{i \in I \setminus J}1) \cdot (\prod_{j \in J} v_j) = \prod_{j \in J} v_j \qedhere \]
\end{proof}

\subsection{The variable graph and the cyclic variables}

\begin{dfn}
  We define the \emph{variable graph} of $E$ to be the directed graph $G$ with $X$ as its
  set of vertices and an edge from $x$ to $y$ whenever $p_x$ contains $y$.
\end{dfn}

\begin{notn}
  Let $\Xcyc$ consist of all elements of $X$ that lie on a cycle in $G$.
\end{notn}

Note that we always have $x \in \downg{x}$, but not always $x \in \Xcyc$.

\begin{expl}
Here is an example of a system $E$ and its variable graph.
\[
\begin{aligned}
  x &= 1 + y^2 + z\\
  y &= y + 1\\
  z &= z + y^2
\end{aligned}
\hspace{3cm}
\begin{tikzcd}
	z & y \\
	x
	\arrow[from=1-1, to=1-1, loop, in=60, out=120, distance=5mm]
	\arrow[from=1-1, to=2-1]
	\arrow[from=1-2, to=1-1]
	\arrow[from=1-2, to=1-2, loop, in=60, out=120, distance=5mm]
	\arrow[from=1-2, to=2-1]
\end{tikzcd}\]
Note the above system has no nullable variables: the $z$ can generate a 1 via $y$.

Nevertheless, there are polynomials in this system that are not recurrent, namely
$x$ and all the polynomials that contain it.  This is due to $x$ not lying on a cycle.
Such variables lead to violation of Proposition \ref{p:gtE-cov}.
Here $1 + z \gtE x$; in fact, $1 + z =_E x$. But $1 + z$ cannot unfold to anything
containing $x$.
\end{expl}

\begin{dfn} \hfill
  \begin{itemize}
    \item
    For $x \in X$, let $\downg{x}$ be the set of all vertices reachable from $x$ in zero or more steps.
    \item
    For $Y \subseteq X$, we also write $\downg{Y}$ for $\bigcup_{y \in Y} \downg{y}$.
    \item
      $V \sse X$ is a \emph{basis} for $G$ if $\downg{V} = X$.
  \end{itemize}
\end{dfn}

\begin{lemma} \label{l:path-covers}
    For $x,z \in X$, $z \in \downg{x}$ iff $x \cov z$.
\end{lemma}
\begin{proof}
$(\RA)$
We proceed by induction on the path from $x$ to $z$.

For the base case, trivially $x \cov x$.

For the inductive case, suppose the shortest path from $x \to z$ looks like 
\[ x \to \cdots \to y \to z\]
By induction hypothesis, $x \cov y$.

Now $y \lra p_y = z \cdot m + r$, so $y \cov z \cdot m$.

By Lemma \ref{l:sse-mon-cov}, $z \cdot m \cov z$. 

Now $x \cov y \cov z \cdot m \cov z$. 

$(\LA)$ Suppose $x \cov z$.

We show $z \in \downg{x}$ by induction on the length of the shortest reduction $x \thra z + r$.

The base case is obtained when $z = x$, in which case $z \in \downg{x}$ by a path of length zero.

So suppose $x \lra p_x \thra z + r$. 

By Proposition \ref{p:cov1}, there exists a monomial $m$ in $p_x$ such that $m \thra z + r'$.
Note that the length of this reduction is yet shorter than that of $x \thra z + r$.

As per Proposition \ref{p:cov2}, we have $m \thra z + r'$ implies $y \thra z + r''$ for some $y$ in $m$.
This is a again a shorter reduction (being part of the original one), so the induction hypothesis applies.

So $z \in \downg{y}$.

Since $y$ occurs in $p_x$, there is an edge from $x$ to $y$ in $G$ and $z \in \downg{x}$.
\end{proof}



\subsection{Cyclic variables}
It would therefore help our analysis to restrict attention to the subsystem of
$E$ consisting only of variables lying on a cycle in $G$. That there is no loss of generality
resulting from this restriction can be seen as follows.

\newcommand{\Ecyc}{E_\circlearrowleft}
\begin{notn}
  \begin{itemize}
    \item $\Xacyc = \setof{x \mid x \notin \Xcyc}$
    \item $\sigma_1 : X \to \nat[X]$ is given by
    \[\sigma_1(x) = \begin{cases}
          x & x \in \Xcyc\\
          p_x & x \in \Xacyc
        \end{cases} \]
      \item Given any function $\sigma : V \to \nat[X]$, let
      $\tilde\sigma : \nat[V] \to \nat[X]$ be the rig homomorphism
      induced by the universal property of the free rig $\nat[V]$.
      % \item $\sigma_E : X \to \nat[X]$ is given by $\sigma_E(x) = p_x$
      \item Given any function $f : A \to A$, let $f^n : A \to A$ be
      $n$-fold composition of $f$ with itself.
    \item $\sigma^* : X \to \nat[\Xcyc]$ be given by
      \[ \sigma^*(x) = \tilde\sigma_1^{|X|}(x) \]
    \item Let $\Ecyc$ be obtained from $E$ by removing all equations
    whose left sides are in $\Xacyc$ and applying $\sigma^*$ to the
    right sides of all equations that remain.
  \end{itemize}
\end{notn}

\begin{prop}\hfill
  \begin{enumerate}
    \item For any $p \in \nat[X], p =_E \sigma_1(p)$.
    \item For any $p \in \nat[X], p =_E \sigma^*(p)$.
  \end{enumerate}
\end{prop}
\begin{proof}
  \begin{enumerate}
    \item Immediate.
    \item By induction on $n$, using the previous part,
    we can show $\tilde\sigma_1^n(p) = p$.

    The claim follows by taking $n = |X|$. \qedhere
  \end{enumerate}
\end{proof}

\begin{prop}
  $\nat[X]/E \simeq \nat[\Xcyc]/\Ecyc$.
\end{prop}
\begin{proof}
  % By the previous proposition, every equation in $\Ecyc$ is validated by $E$.
  % \greyout{It follows that the inclusion $\iota : \Xcyc \sse X$ and $\sigma^*$ induce
  % homomorphisms on the quotients that are inverses to each other.}
  %
  By the previous proposition, the $E$-equivalence class of every $x \in \Xacyc$ has a
  representative in $\nat[\Xcyc]$, namely $\sigma^*(x)$.  This yields the inverse
  to the homomorphism induced by the inclusion $\iota : \Xcyc \sse X$.
\end{proof}
\begin{proof}[More detail.]
  The rig homomorphism $\tilde\iota : \nat[\Xcyc] \to \nat[X]$
  induced by the inclusion $\iota : \Xcyc \subset X$ maps
  any polynomial in $\nat[\Xcyc]$ to itself,
  and in particular any equation $x = p'_x$ in $\Ecyc$
  the itself as well.  All such equations are validated by $E$
  since $p_x =_E p'_x$ by the above proposition.

  $\tilde\iota$ thus induces a rig homomorphism $\hat \iota :
  \nat[\Xcyc]/\Ecyc \to \nat[X]/E$.

  The inverse of $\hat \iota$ is given by $\hat \sigma^*$, which in turn
  is induced by $\tilde \sigma^*$ because $\sigma^*$ maps every equation
  in $E$ to $\Ecyc$-equal terms in $\nat[\Xcyc]$.

  That $\hat \sigma^*(\hat \iota (x)) = x$ follows by the proposition,
  whereas $\hat \iota (\hat \sigma^*(x)) = x$ follows for the same
  reason when $x \in \Xcyc$ and when $x \notin \Xcyc$,
  $x$ gets mapped to $\sigma^*(x) \in \nat[\Xcyc]$ which is $E$-equal to $x$.
\end{proof}

In light of the proposition above, there is no loss of generality
in the following.

\begin{asmn}
  \label{asmn:cyc}
   \fbox{$X = \Xcyc$}
 \end{asmn}

The above assumption is hereby admitted for the rest of the paper.

The key result of this section is that, within the scope of the assumptions above,
every variable is recurrent.

\begin{prop} \label{p:cycvar}
Under assumptions \ref{asmn:null} and \ref{asmn:cyc}, $X = \Xsc$.
\end{prop}
\begin{proof}
Let $x \in X$. We will show that $x$ is recurrent.

By Corollary \ref{p:sc=cov}, it suffices to show $p_x \cov x$.

Pick a cycle $C$ containing $x$, say 
\[ x \to y \to \cdots \to x\]

By Lemma \ref{l:path-covers}, $y \cov x$.

Also $p_x = y \cdot m + r$, so $p_x \cov y \cdot m$.

By Lemma \ref{l:sse-mon-cov}, $y \cdot m \cov y$, so $p_x \cov y \cov x$.

Indeed, $x \in \Xsc$.
\end{proof}

\begin{cor}\label{c:recur}
  Under assumptions \ref{asmn:null} and \ref{asmn:cyc}, every polynomial is recurrent.
\end{cor}
\begin{proof}
  By Proposition \ref{p:cycvar} and Corollary \ref{p:sc=cov}.3.
\end{proof}


\subsection{High polynomials and the basis of the variable graph}

The assumptions are assumed throughout this section.

Under these assumptions, we have $X = \Xsc$, all polynomials are recurrent, and $\gtE = \cov$ 
by Propositions \ref{p:cycvar} and \ref{p:gtE-cov}.
We will this mix these notations throughout as convenient.


\begin{dfn}
  $p \in \nat[X]$ is \emph{high} iff $p \gtE q$ for all $q \in \nat[X]$.
\end{dfn}

% \begin{lemma} \hfill
%   \begin{enumerate}
%     \item If $m \neq 1$ is a monomial, then $\const{n} \not \gtE m$.
%     \item For any non-constant monomial $m$, $m \gtE k \cdot m$.
%     \item For any non-constant polynomial $p$, $p \gtE k \cdot p$.
%   \end{enumerate}
% \end{lemma}
% \begin{proof}
%   \begin{enumerate}
%     \item
%     By Proposition \ref{p:gtE-cov} and the fact that $\const{n}$ is a normal form for unfolding reductions.
%     \item Note that
%     \item
%   \end{enumerate}
% \end{proof}



\begin{lemma} \label{l:high} \hfill
\begin{enumerate} 
    \item \label{l:high1} 
    For any variable $x$, and monomials $m_1$, \dots $m_k$, $x \cdot \prod m_i \cov \sum m_i$.
    \item \label{l:high2}
    $p \in \nat[X]$ is high iff $p \gtE m$ for all monomials $m$.
    \item \label{l:high3}
    If a high polynomial exists, then a high monomial exists.
\end{enumerate}
\end{lemma}
\begin{proof}
\begin{enumerate}
    \item 
    Let $m_1,\dots,m_k$ be given and let $x \in X$.

    Put $h = \prod m_i$. By Lemma \ref{l:sse-mon-cov}, $h \cov m_i$ for each $i$.

    By Proposition \ref{p:xcovn}, $x \cov k$.  So 
    \[ x \cdot h \cov k \cdot h = \sum_{i=1}^k h \cov \sum_{i=1}^k s_i \]
    \item  
      $(\RA)$ Trivial.

  $(\LA)$ Suppose $p \cov m$ for all monomials $m$.
  
  Let $q = \sum m_i$ be given.  By hypothesis, $p \cov x \cdot \prod m_i$.
  
    By part \ref{l:high1}, $x \cdot \prod m_i \cov q$.  So $p \cov q$.
  
    \item
    Let $p = \sum_{i \in I} m_i$ be a high polynomial composed of monomials $m_i$.
    
    By part \ref{l:high1}, $x \cdot \prod m_i \cov p$, where $x$ is any variable.
    
    Since $p$ is high, so is the monomial $x \cdot \prod m_i$. \qedhere
    \end{enumerate}
\end{proof}

\begin{thm} \label{t:high}
  \begin{enumerate}
    \item \label{t:high1} A polynomial is high iff it contains a high monomial.
    \item \label{t:high2} A monomial $m = \prod_{v \in V} v$ is high iff $V$ contains a basis of compactible elements.
    \item \label{t:high3} If $\nat[X]/E$ contains high elements iff such a basis exists.
  \end{enumerate}
\end{thm}

\begin{proof}
  \begin{enumerate}

    \item
    $(\LA)$
    Suppose $p = m + q$ and $m$ is high.

    Let $p'$ be given. Since $m$ is high, $m \gtE p'$.

    Let $r$ be such that $m =_E p' + r$.

    Then $p =_E p' + r + q$ and $p \gtE p'$.
    
    $(\RA)$
    Let $p = \sum_{i \in I} s_i$ be a high polynomial, where each $s_i$ is a monomial.  By Lemma \ref{l:high}.\ref{l:high3}, a high monomial $m$ exists.  Because $p$ is high,
    \[ p \gtE m \]
    By Lemma \ref{p:gtE-cov} together with Corollary \ref{c:recur}
    \[ p \cov m \]
    We thus have that there exists a polynomial $r$ such that
    \[ p \thra m + r\]
    By Proposition \ref{p:sum-red}, there exists a monomial $s_j$ in $p$ and a polynomial $r'$ such that $s_j \thra m + r'$.  We thus have that $s_j \cov m$, so $s_j \gtE m$.
    Since $m$ is high, and $s_j \gtE m$, $s_j$ is high and $p$ contains a high monomial.

    \item 
    $(\LA)$
    Suppose $m_0 = m_0' \cdot \prod {v_i}$, and $B = \setof{v_1,\dots,v_k}$ is a basis for the graph.

    By Lemma \ref{l:high}.\ref{l:high2}, it suffices to show that $m_0 \cov m$ for any monomial $m$.

    Let $m = \prod_{j = 1}^l x_j$.

    Since $B$ is a basis, $x_j \in \downg{B}$.  By Lemma \ref{l:path-covers} and \ref{l:sse-mon-cov}, 
        $b \cov x_j$ for each $j$.

    Then $b^l \cov m = \prod x_j$.

    By Proposition \ref{p:compact-n}, $v_i \cov v_i^l$ for each $v_i \in B$.  So $b \cov b^l$.

    Now $b \cov b^l \cov \prod x_j = m$. 
    
    By Lemma \ref{l:sse-mon-cov}, $m_0 \cov b$. Hence $m_0 \cov m$.

    $(\RA)$ 
    Let $m = \prod_{i = 1}^k v_i$ be a high monomial.

    Let $V_0 = \setof{v_i | 1 \le i \le k}$ be the set of variables in $m$.

    We define a function $f : V_0 \to V_0$ as follows. 
    Given $v \in V_0$, note that $m \cov v^{k+1}$, as $m$ is high.
    By Proposition \ref{p:cov2} and the Pigeonhole principle, there must exist a $w \in V_0$ such that $w \cov v^{2}$.
    Let $f(v)$ be defined by choosing such a $w$.

    For $n \ge 0$, let $V_{n+1} = f[V_n]$. Note that $V_{n+1} \sse V_n$.  Since all $V_n$ are finite, 
    there exists $N$ such that $V_{N+1} = V_N$.

    Then $f \rightharpoonup V_N$, the restriction of $f$ to $V_N$, is a bijection.

    We claim that $V_N$ is a basis, and every element of $V_N$ is compactible.  This follows by the following sequence of claims.

    \begin{enumerate}
    \item By construction, $f(v) \cov v^2$
    \item By Lemma \ref{l:sse-mon-cov}, $f(v) \cov v$.
    \item By induction, $f^i(v) \cov v$ for all $i \ge 0$.
    \item When restricted to $V_N$, $f$ is a bijection on a finite set.  For each $v \in V_N$, there thus exists a $l \ge 0$ such that $f^l(f(v)) = v$.
    \item Combining these facts, we get $v = f^l(f(v)) \cov f(v) \cov v^2$ for each $v \in V_N$.  So elements of $V_N$ are compactible.
    \item Moreover, for each $v_i \in V_0$, $f^N(v_i) \cov v_i$. 
    \item Let $u = \prod_{v \in V_N} v$.  Since each $v \in V_N$ is compactible, $u \cov u^k$.
    \item $u \cov v_i$ for all $i$ by combining the previous two items.
    \item So $u \cov u^k \cov m$. Thus $u$ is high.
    \item Now $u \cov x$ for all $x \in X$. 
    \item By Proposition \ref{p:cov2}, $v \cov x$ for some $v \in V_N$.
    \item Now $x \in \downg{V_N}$ by Proposition \hired{still need to prove it}.
    \end{enumerate}        
    
    \item $(\RA)$ Suppose a high polynomial exists in $\nat[X]/E$.
    By part \ref{t:high1}, it contains a high monomial.
    By part \ref{t:high2}, the latter contains a basis of compactible elements.
    Hence such a basis exists.

    $(\LA)$. Assume such a basis exists, say, $B = \setof{v_1, \dots, v_n}$.
    By part \ref{t:high2}, $\prod v_i$ is high.  Hence, high elements exist.    
  \end{enumerate}
\end{proof}

\begin{thm}
  Assume Assumptions \ref{asmn:null} and \ref{asmn:cyc}.

  High elements exist in $\nat[X]/E$ exist iff $G$ contains a basis of compactible elements.
\end{thm}

\newpage 
\section{Examples}
\begin{enumerate}
    \item 
\[ 
\begin{cases}
    x = x + 1\\
    y = y^2 + x
\end{cases}
\]
    \item 
\[ 
\begin{cases}
    x = 2x + 1\\
    y = y^2 + x
\end{cases}
\]
Here we can prove that $Y^2 = Y + 1$, as follows.
\begin{align*}
y = y^2+x=y^2+2x+1\\
y^2+x = y + x + 1 \\ 
y+1 = y^2 + x +1 = y + x + 2\\ 
y = y^2 + 2x + 1 \\ 
y^2 = y^3+yx 
= y^3 + 2xy + y 
= y^3+2yx+y^2+2x+1 \\ 
= y^2+2yx+y+(x+1) \\ 
=y^3+xy + (x+1) \\ 
y^2+x+1
\end{align*}
\end{enumerate}
\newpage
\section{OLD! Characterizing the high elements}
\begin{enumerate}
  \item Eliminate ``zero variables": those that do not contain a constant term
  after $n$ iterations of $\bar p$.
  \item Generate the digraph of inclusion between variables: an edge from $i$ to $j$ if
  $p_j$ contains $x_i$.
  \item Remove everything not on a cycle: these can be defined in terms of what's left.
  \item Consider the terminal SCCs: the high elements necessarily contain one from each.
  \item If each satisfies $x^2 \le x$ then this is also sufficient, so
  the high elements are exactly those that contain a monomial containing a variable
  from each terminal SCC.
  \item So this is the condition: each TSCC contains an $i$ such that $p_i$ contains
  a monomial with two variables from $\cc{i}$.
\end{enumerate}

More comments:
\begin{itemize}
  \item Given $E = \setof{\bar x = \bar p(\bar x)}$, we split it into
  $E = E_0 \cup E_1 \cup E^*$, where
  $E_0$ contains all variables that do not contain a constant,
  $E_1$ contains all variables that do not lie on a cycle,
  and $E^*$ is the ``core" that the analysis will be performed on.
  \item To define $E_0$, can either use induction or the following.

  We denote by $\bar p$ the operator on $\nat[\bar x]^n$ defined tautologically by
    \[ (\Xacyc,\dots,x_n) \longmapsto (p_1(\bar x),\dots,p_n(\bar x)) \]

  % Consider the following operator $P = P_E : \nat[\bar x]^n \to \nat[\bar x]^n$ given by
  % \[ P(t_1,\dots,t_n) = (p_1(\bar t),\dots,p_n(\bar t)) \]

  We consider any element $t \in \nat[\bar x]$ in its sum-of-monomials normal form, so that the degree,
  any constant terms, etc. are immediate given the polynomial.

  The condition defining elements of $E_0$ now becomes:
  \[ i \in E_0 \iff \bar p^n(\bar 0)_i = 0 \]
  This is equivalent to $\bar p^n(\bar x)_i$ does not contain a constant term.

  \item WLOG, we assume $E = E^*$, and so keep calling it $E$.
  \item Note that $x_i \neq n$ for any $n \in \nat$, once $E_1$ is removed,
  since $n$ does not contain any variables and hence $x_i$ cannot lie
  on a cycle.
\end{itemize}

From now on, assume $\bar x = \bar p (\bar x)$.
\textsc{Proposition.} Assume $E = E^*$, we have
\begin{enumerate}
  \item $x_i \neq n$, if $n \in \nat$.
  \item For any non-zero monomial $m \in \nat[\bar x]$, $1 \le m$.
  \item For any $x_i$ and non-zero $m$, $x_i \le x_i \cdot m$.
  \item $\le$ is a congruence wrt multiplication.
\end{enumerate}

\textsc{Lemma.} Given $x_i \in \bar x$, TFA:
\begin{enumerate}
\item $i \notin E_0$
\item $\bar p^n(\bar x)_i$ contains a constant term.
\item $1 \le x_i$.
\end{enumerate}

\begin{proof}
  \begin{description}
  \item[$(1) \RA (2)$.] We proof the contrapositive.

  Suppose $\bar p^n(\bar x)_i$ does not contain a constant term, so that
  \[\bar p^n(\bar x)_i = \sum_{k \in K} x_{j(k)}\cdot m_k \] where
  $m_k$ is a non-zero monomial and $x_{j(k)} \in \bar x$.

  Now $\bar p^n(\bar 0)_i = \sum_{k \in K} 0 \cdot m_k[\bar 0/\bar x] = 0$.

  That is $i \in E_0$.

  \item[$(2) \RA (3)$.] Suppose $\bar p^n(\bar x)_i$ contains a constant term.

  Then $\bar x = \bar p(\bar x)$ implies $\bar x = \bar p^n(\bar x)$.

  Now $x_i = \bar p^n(\bar x)_i = q + 1$ for some $q \in \nat[\bar x]$ by assumption.

  \item[$(3) \RA (1)$.]

  Let $q \in \nat[\bar x]$ and suppose $E \models 1 + q = x_i$.

  Then $\semof{1 + q}_\rho = \semof{x_i}_\rho$ in $\nat[\bar x]/{E^*_{\bar x}}$ for all $\rho$.

  In particular, take $\rho(x_i) = 0$.  Note that $\semof{q}_\rho \in \nat$, hence non-negative.

  Thus $\semof{1 + q} = 1 + \semof{q} > 0$.

  On the other hand, since $E \models x_i = \bar p^n(\bar x)_i$, we have
  \[  \semof{\bar p^n(\bar x)_i}_{\bar 0} = \bar p^n(\bar 0)_i = \semof{x_i}_{\bar 0} > 0 \]
  whence $i \notin E_0$.
\end{description}

\end{proof}

From now on, fix $E = E^*$.
\textsc{Lemma.}

For all $x_i \in \bar x$, $2 \le x_i$.

\begin{proof}
By hypothesis, $1 \le x_i$ for all $i$, hence $1 \le m$ for any non-constant monomial in $\nat[\bar x]$.

Write $x_i = p_i = q + 1$, where $q$ contains a non-constant monomial because $x_i$ is non-constant.

So $1 \le q$, $q = q' + 1$, $x_i = q' + 2$, and thus $2 \le x_i$.
% Note that $q = x_i + q'$ because $x_i$ lies on a cycle.
% Now $x_i = q + 1 = q' + x_i + 1 = q' + q' + x_i + 1 + 1$, whence $2 \le x_i$.
\end{proof}

\textsc{Lemma.}


From now on, let $G = G(E)$ be the variable inclusion digraph, where
$i$ has an edge to $j$ iff $p_j$ contains $x_i$.

\textsc{Lemma.} If there's a path from $i$ to $j$ then $x_i \le x_j$.

\begin{proof} Induction on the path.
  \begin{description}
    \item[Base case] $x_i \le x_i$ is trivial.
    \item[Induction] Suppose $p_j$ contains $x_i$, and there's a path from
    $j$ to $k$.

    By induction hypothesis, we have $x_j \le x_k$.

    By transitivity, it suffices to show that $x_i \le x_j$.

    By hypothesis $p_j = x_i \cdot m + q$ for some non-zero $m$, so
    $x_i \le x_i \cdot m \le p_j = x_j$.
  \end{description}
\end{proof}

\textsc{Lemma.} Converse to the above.
\begin{proof}
  \textcolor{red}{NEED TO DO THIS!!!}
\end{proof}

\textsc{Definition.} A variable $x_i \in \bar x$ is called a
\emph{terminal} if it is maximal with respect to $\le$ restricted to variables:
\[ \forall j.  x_i \le x_j \implies x_j \le x_i \]

\newcommand{\term}{\mathcal{T}}

\textsc{Definition.}
\begin{enumerate}
  \item $\uparrow p = \setof{q \mid p \le q}$
  \item $\downarrow p = \setof{q \mid q \le p}$
  % \item $\cc{p} = {\uparrow p} \cap {\downarrow p} = \setof{q \mid p \le q, q \le p}$.
  \item $H = H(E) = \setof{p \mid \downarrow p = \nat[\bar x]}$
  \item $p$ is \emph{high} if $p \in H$.
  \item $T = \setof{x_i \mid x_i \text{ is terminal }}$
  \item $\cc{x_i} =  \setof{x_j \mid x_i \le x_j, x_j \le x_i}$.
  \item A \emph{basis} for a graph is a set of vertices $B$ such that,
  every node can reach one of the vertices in $B$.
  \item A \emph{minimal basis} is a basis with least cardinality.
  \item An element satisfying $p^2 \le p$ is called \emph{compactible}.
\end{enumerate}

\textsc{Proposition.}
\begin{enumerate}
  \item $p \in H, p \le q \then q \in H$
  \item $x \in T$, then $y \in \uparrow x \iff \cc{y} = \cc{x}$
\end{enumerate}
\textcolor{red}{DO THIS}

\textsc{Proposition.} For $x_i \in T$, TFAE:
\begin{enumerate}
  \item $x_i^2 \le x_i$
  \item $\exists x_j \in \cc{x_i}$ s.t.
  $p_j(\bar x) = m(\bar x) + q(\bar x)$, where $m$ contains two elements of $\cc{x_i}$.
  \item $\exists x_j \in \cc{x_i}. \deg_{x_j} p_j(\bar v) \ge 2$, where
  \[ v_k = \begin{cases}
  x_j &x_k \in \cc{x_i} \\
  1 &o/w
\end{cases} \]
\end{enumerate}

\begin{proof}
\begin{description}
  \item[(1) $\RA$ (2)]
  Suppose $x_i \simeq x_i^2 + q$.

  Since $x_i$ lies on a cycle, we can write $x_i \simeq x_i + q'$, where $q' \neq 0$.

  \item[(2) $\RA$ (1)] We have
  \begin{align*}
    x_j \simeq p_j(\bar x) &= m(\bar x)+q(\bar x) \\
    &\ge m(\bar x)\\
    &= x_{k_1} \cdot x_{k_2} \cdot m'(\bar x) && (x_{k_1},x_{k_2} \in \cc{x_i}, m' \neq 0)\\
    &\ge x_i \cdot x_i \cdot 1\\
    &= x_i^2
  \end{align*}
\end{description}
\end{proof}

\textsc{Theorem.} TFAE
\begin{enumerate}
  \item
  $\nat[\bar x]$ has high elements modulo $E$
  \item
  $G(E)$ contains a minimal basis of compactible elements
  \item
  The terminal elements of any basis for $G(E)$ are compatible.
\end{enumerate}

\begin{proof}[Proof sketch.]
  Every compactible terminal $x_i$ is high for the subgraph that it generates.

  The high elements for all of $G(E)$ are those polynomials that contain a
  monomial that contain a variable from each terminal class.
\end{proof}

\section{Comparing systems of equations}
Let $E, E'$ be given.
\begin{enumerate}
  \item (E) Alpha conversion (useless)
  \item (E) Syntactic equivalence:
  $E'$ can result from $E$ by rewriting RHS using $E$. (For systems in the same set of variables.)
  \item (E) Mu rule equality (equivalent to the above?)
  \item (E) $E'$ results from a sequence of ``adding or removing variables'' of $E$:
  \begin{itemize}
    \item If $x$ does not lie on a cycle, replace
    $E$ with $E\setminus(x = p_x(\bar x))[x := p_x(\bar x)]$
    \item For any term $t(\bar x)$, replace $E$ with
    $E \cup (y = t(\bar x))$
  \end{itemize}
  \item (E) $\nat[V]/E \simeq \nat[W]/E'$
  \item (E) Equivalent to the above? Same but $\nat[V]/E$, $\nat[W]/E'$ are categories (groupoids)
  and the equivalence is groupoid equivalence.
  \item (I) $E \models E'$ if for any semiring $S$
  $E \models_{Rig} S \implies E' \models_{Rig} S$
  \item (I) $E(V) \preceq E'(W)$ if $\exists \sigma : V \to \nat[W]$ s.t.
   \[ \forall (l = r) \in E, \qquad \nat[W]/E' \models l[\sigma] = r[\sigma] \]
  \item (I) If $S$ admits a solution to $E$, then $S$ must admit a solution to $E'$.
\end{enumerate}

\textsc{Conjecture.}
If $h : \nat[\bar x]/E \to \nat[\bar x]/E$ is an injective endomorphism, then it's an automorphism.

\end{document}
